\chapter{Recovery layers and variational limits}\label{sec:phase-recovery}
The lower bound becomes sharp by placing the one-dimensional profile along an interface. We develop the construction behind \cite{ChAC}, Proposition 3.6, first for smooth interfaces and then for compactly contained finite-perimeter sets. The general relative-boundary approximation is recorded separately as an external input. Throughout, $U\subset\mathbb R^n$ is bounded and open and $n\geq2$. We write $q(t)=\tanh(t/\sqrt2)$ for the profile of Example \ref{ex:allen-cahn-layer}.

\begin{proposition}[A flat layer and its limiting energy measure]\label{prop:phase-flat-interface}
Let $e\in\mathbb S^{n-1}$, $b\in\mathbb R$, and $u_\varepsilon(x)=q((x\cdot e-b)/\varepsilon)$ on $\mathbb R^n$, where $q(t)=\tanh(t/\sqrt2)$. This is a smooth Allen--Cahn solution, and locally
\[
 u_\varepsilon\to2\chi_{\{x\cdot e>b\}}-1,\qquad
 \mu_\varepsilon^{u_\varepsilon}\rightharpoonup
 c_W\mathcal H^{n-1}\llcorner\{x\cdot e=b\}.
\]
The phase convergence is in every finite $L^p$ on bounded sets.
\end{proposition}
\begin{proof}
The chain rule and $q''=W'(q)$ give the equation. The phase convergence follows pointwise off the hyperplane and then by bounded convergence. Write $x=y+(b+\varepsilon t)e$ with $y\in e^\perp$. Equipartition gives energy density $\varepsilon^{-1}(q'(t))^2$. For a continuous compactly supported test function $\zeta$, Fubini gives
\[
 \int\zeta\,d\mu_\varepsilon
 =\int_{e^\perp}\int_{\mathbb R}
 \zeta(y+(b+\varepsilon t)e)(q'(t))^2\,dt\,d\mathcal H^{n-1}(y).
\]
Its integrand is dominated by an integrable function using the bounded transverse projection of the support and $\int(q')^2=c_W$. Dominated convergence proves the measure limit. The fixed translation belongs inside the scaled argument: $q((x\cdot e-b)/\varepsilon)$ keeps the interface at $x\cdot e=b$. This clarifies the rescaling in \cite{ChAC}, Section 2.3.1.
\end{proof}

\begin{lemma}[Tubular coordinates near a smooth boundary]\label{lem:phase-tubular-coordinates}
Let $E$ be open with compact $C^2$ boundary $\Sigma=\partial E\Subset U$. For some $\rho>0$, the map
$F(y,t)=y+t\nu(y)$, $y\in\Sigma$, $|t|<\rho$, is one-to-one onto a tubular neighborhood, where $\nu$ is outward. Its volume element is $J(y,t)d\mathcal H^{n-1}(y)dt$, with
\[
 J(y,t)=\det(I+tD_\Sigma\nu),\qquad J(y,0)=1,
\]
and $J>0$ and $J\to1$ uniformly as $t\to0$. The signed distance $d$, positive in $E$, satisfies $d(F(y,t))=-t$, $|Dd|=1$, and $\Delta d|_\Sigma=-H$ for $H=\operatorname{div}_\Sigma\nu$.
\end{lemma}
\begin{proof}
At $t=0$, the derivative of $F$ maps tangent directions identically and the $t$ direction to $\nu$, so the inverse function theorem gives local coordinates. Compactness supplies a uniform radius. Global injectivity for a smaller radius follows by contradiction: two colliding normal segments with lengths tending to zero would have base points tending to the same boundary point, contradicting local injectivity. Nearest-point minimization shows these normal coordinates give the signed distance. Tangential differentiation of $F$ yields $I+tD_\Sigma\nu$, and its normal derivative is $\nu$; this gives the Jacobian. Positivity and uniform convergence follow by reducing $\rho$. Since $Dd=-\nu$ on the boundary and has zero normal derivative along normal segments, its divergence there is $-\operatorname{div}_\Sigma\nu$. These are the coordinate facts used in \cite{ChAC}, p.~10.
\end{proof}

\begin{theorem}[Recovery along a smooth compact interface]\label{thm:phase-smooth-recovery}
Suppose $E\Subset U$ is open with $C^2$ boundary. For nonempty $E$, let $d$ be its signed distance, positive inside, and put $u_\varepsilon=q(d/\varepsilon)$ with $q(t)=\tanh(t/\sqrt2)$; for empty $E$, put $u_\varepsilon=-1$. Then $u_\varepsilon\in H^1(U)\cap L^4(U)$, $|u_\varepsilon|\leq1$, and
\[
 u_\varepsilon\to u_E\text{ in every finite }L^p(U),\qquad
 E_\varepsilon(u_\varepsilon;U)\to c_WP(E;U).
\]
Moreover the energy measures converge to $c_W\mathcal H^{n-1}\llcorner\partial E$ on compactly supported continuous tests.
\end{theorem}
\begin{proof}
Signed distance is Lipschitz with constant one, so the composition is Sobolev and bounded. It converges to the two phases off the measure-zero boundary. In the tube of Lemma \ref{lem:phase-tubular-coordinates}, $|Dd|=1$ and the energy is
\[
 \int_\Sigma\int_{-\rho/\varepsilon}^{\rho/\varepsilon}
 (q'(s))^2J(y,-\varepsilon s)\,ds\,d\mathcal H^{n-1}(y).
\]
The bounded Jacobian and integrability of $(q')^2$ allow dominated convergence, giving $c_W\mathcal H^{n-1}(\Sigma)$. Outside the tube, $|d|\geq\rho$. Both $(q'(d/\varepsilon))^2$ and $W(q(d/\varepsilon))$ are bounded by $Ce^{-c\rho/\varepsilon}$; the gradient bound $|Dd|\leq1$ therefore makes the remaining energy at most $C|U|\varepsilon^{-1}e^{-c\rho/\varepsilon}\to0$. Insert a test function in the same calculation to obtain the measure limit. Bounded convergence proves all the stated phase convergences. This is the smooth-interface construction in \cite{ChAC}, Proposition 3.6, with the tail estimate retained. Empty $E$ is handled by $u_\varepsilon\equiv-1$.
\end{proof}

\begin{proposition}[The equation error of a curved recovery layer]\label{prop:phase-recovery-residual}
In the tubular neighborhood of the preceding theorem,
\[
 \varepsilon^2\Delta u_\varepsilon-W'(u_\varepsilon)
 =\varepsilon q'(d/\varepsilon)\Delta d.
\]
On the interface this equals $-\varepsilon q'(0)H$. Thus a recovery sequence need not consist of critical points; even when $H=0$ on the central interface, the displayed error need not vanish away from it.
\end{proposition}
\begin{proof}
Differentiate the composition twice to get
$\varepsilon^2\Delta u_\varepsilon=q''(d/\varepsilon)|Dd|^2+\varepsilon q'(d/\varepsilon)\Delta d$. Use $|Dd|=1$, $q''=W'(q)$, and the boundary value of $\Delta d$ from Lemma \ref{lem:phase-tubular-coordinates}. This explains the distinction between recovery competitors in \cite{ChAC}, Section 3.1, and constructed solutions in Section 4.
\end{proof}

\begin{theorem}[Recovery for a compactly contained finite-perimeter set]\label{thm:phase-interior-recovery}
Let $E$ agree almost everywhere with a subset of a compact set $K\Subset U$, and assume $P(E;\mathbb R^n)<\infty$. There are $u_\varepsilon\in H^1(U)\cap L^4(U)$, $|u_\varepsilon|\leq1$, such that
\[
 u_\varepsilon\to u_E\text{ in }L^1(U),\qquad
 E_\varepsilon(u_\varepsilon;U)\to c_WP(E;U).
\]
\end{theorem}
\begin{proof}
Use the mollification and regular-level construction in Theorem \ref{thm:perimeter-smooth-approximation}. Since the original set is supported in $K$, the approximating sets $E_k$ lie in shrinking neighborhoods of $K$; hence $E_k\Subset U$ for all sufficiently large $k$. Their volumes converge in symmetric difference and their total perimeters converge to $P(E;\mathbb R^n)=P(E;U)$.

For each $k$, Theorem \ref{thm:phase-smooth-recovery} supplies a family $v_{k,\varepsilon}$. Choose numbers $\eta_k\downarrow0$ so that for $0<\varepsilon<\eta_k$ its $L^1$ error from $u_{E_k}$ and its energy error from $c_WP(E_k;U)$ are each at most $1/k$. Choose $k(\varepsilon)\to\infty$ with $\varepsilon<\eta_{k(\varepsilon)}$ and set $u_\varepsilon=v_{k(\varepsilon),\varepsilon}$. The triangle inequality, the factor-two phase distance, and convergence of the perimeters prove the result. The empty approximation is assigned the constant minus-one family. This proves the interior case of \cite{ChAC}, Proposition 3.6 from the BV approximation already available here.
\end{proof}

\begin{definition}[Gamma convergence in the phase topology]\label{def:phase-gamma-convergence}
Functionals $F_\varepsilon:L^1(U)\to[0,\infty]$ Gamma-converge to $F_0$ if (i) every sequence $\varepsilon_j\downarrow0$, $u_j\to u$ in $L^1$ satisfies $F_0(u)\leq\liminf_jF_{\varepsilon_j}(u_j)$, and (ii) for each $u$ there is a family $u_\varepsilon\to u$ in $L^1$ with $\limsup_{\varepsilon\downarrow0}F_\varepsilon(u_\varepsilon)\leq F_0(u)$. The second family is called a recovery sequence. This is the variational convergence described in \cite{ChAC}, Remark 3.7; it concerns energy values and minimization, not convergence of all critical points.
\end{definition}

\begin{theorem}[Relative recovery without prescribed boundary values]\label{thm:phase-relative-recovery}
Let $U$ be bounded with Lipschitz boundary, and let $E\subset U$ have finite perimeter relative to $U$. There exists a recovery family with $|u_\varepsilon|\leq1$,
\[
 u_\varepsilon\in H^1(U)\cap L^4(U),\quad
 u_\varepsilon\to2\chi_E-1\text{ in }L^1(U),\quad
 E_\varepsilon(u_\varepsilon;U)\to c_WP(E;U).
\]
This is \cite{ChAC}, Proposition 3.6, stated on a Lipschitz domain to specify a standard relative approximation setting. The approximation of interfaces meeting $\partial U$, followed by the diagonal recovery construction, remains an external proof dependency. No trace is fixed, and no boundary energy is included in this statement.
\end{theorem}

\begin{corollary}[The unconstrained perimeter Gamma limit]\label{cor:phase-gamma-limit}
On a bounded Lipschitz domain, the energies of Definition \ref{def:phase-local-energy} Gamma-converge in $L^1(U)$ to
\[
 F_0(u)=
 \begin{cases}
 c_WP(E;U),&u=2\chi_E-1\text{ a.e.},\ E\text{ of finite perimeter},\\
 +\infty,&\text{otherwise}.
 \end{cases}
\]
\end{corollary}
\begin{proof}
Theorem \ref{thm:phase-liminf} is the lower-bound condition. For finite $F_0$, Theorem \ref{thm:phase-relative-recovery} is the upper-bound condition. If $F_0(u)=+\infty$, the upper-bound condition is vacuous, and the constant family $u_\varepsilon=u$ in the extended functional domain suffices. This proves the assertion using the explicit external recovery input; it is not a proof of that approximation theorem. Fixed traces or volume constraints would define different admissible problems and require their own recovery arguments.
\end{proof}

\begin{proposition}[Limits of minimizers with a fixed local radius]\label{prop:phase-local-minimizers}
Let $U$ be a bounded Lipschitz domain and $\delta>0$. Suppose $\varepsilon_j\downarrow0$, $E_{\varepsilon_j}(u_j;U)\leq C$, and
\[
 E_{\varepsilon_j}(u_j;U)\leq E_{\varepsilon_j}(v;U)
 \quad\text{whenever }\|v-u_j\|_{L^1(U)}<\delta.
\]
After a subsequence, $u_j\to u_E$ in $L^1(U)$, the energies converge to $c_WP(E;U)$, and
\[
 P(E;U)\leq P(F;U)\quad\text{if }|E\triangle F|<\delta/2.
\]
Only non-strict local minimality is concluded, even if the approximating minima were strict. The radius $\delta$ is independent of $j$.
\end{proposition}
\begin{proof}
Compactness gives the phase limit. For a finite-perimeter competitor $F$ with $2|E\triangle F|<\delta$, choose its recovery family $v_j$ from Theorem \ref{thm:phase-relative-recovery}. Then
$\|v_j-u_j\|_1\to\|u_F-u_E\|_1=2|E\triangle F|<\delta$, so for all large $j$ the comparison is admissible. There is a fixed positive margin below $\delta$. The lower bound for $u_j$, the minimizing inequality, and recovery for $F$ yield
$c_WP(E;U)\leq\liminf_jE_{\varepsilon_j}(u_j;U)\leq\limsup_jE_{\varepsilon_j}(u_j;U)\leq c_WP(F;U)$.
Taking $F=E$ proves energy convergence; arbitrary $F$ proves local perimeter minimality. Infinite-perimeter competitors impose no restriction. This is the argument of \cite{ChAC}, Proposition 3.10, with the phase/indicator distance and strict comparison margin explicit, in the unconstrained Euclidean-domain setting.
\end{proof}
